% GENERATED FILE. Edit canonical lessons, then run npm run generate:books.
% canonical-source-hash: fnv1a64:8a8f03624431fbf2
% canonical-lessons: KA-C232-kharcu, KA-C232-vime, KA-C232-salu, KA-C232-anubhava, KA-C232-rajiname

\chapter{Expense, Insurance, Queue, Experience, Resignation}
\label{ch:ka-expense-insurance-queue-experience-resignation}
\hlchaptermodality{232}

\begin{chapteropening}
\textbf{By the end of this chapter:} \emph{I can say an expense, insurance, a queue, experience and a resignation.}

Complete the last lesson of chapter 232: I can say an expense, insurance, a queue, experience and a resignation.
\end{chapteropening}

\section[\texorpdfstring{\kn{ಖರ್ಚು} (kharcu)}{ಖರ್ಚು (kharcu)}]{\kn{ಖರ್ಚು} (kharcu) — an expense, spending}

\label{lesson:KA-C232-kharcu}



\begin{quote}
Before the new one: say the Kannada for a clerk, then the Kannada for an officer.
\end{quote}

\subsection*{You'll want to know: \kn{ಖರ್ಚು}}
\textbf{\kn{ಖರ್ಚು}} — \emph{kharcu} — \textquotedblleft{}an expense, spending\textquotedblright{}.

\begin{grammarlens}[title={What you've built}]
One more word for this chapter.
\end{grammarlens}

\subsection*{Your turn}
\begin{itemize}
\raggedright
  \item \emph{Say it:} \emph{kharcu}
  \item \emph{Say it:} \emph{kharcu}, once more
  \item \emph{Say it:} say \emph{kharcu}
  \item \emph{Recall:} say the Kannada for a clerk, then the Kannada for an officer, then say \emph{kharcu} again
\end{itemize}

\begin{tcolorbox}[breakable,colback=teal!4,colframe=teal!35!black,title={Before you move on}]
What does \kn{ಖರ್ಚು} mean? (\textquotedblleft{}An expense, spending\textquotedblright{}.) Say it once more.
\end{tcolorbox}
\section[\texorpdfstring{\kn{ವಿಮೆ} (vime)}{ವಿಮೆ (vime)}]{\kn{ವಿಮೆ} (vime) — insurance}

\label{lesson:KA-C232-vime}



\begin{quote}
Before the new one: say the Kannada for an officer, then the Kannada for an expense.
\end{quote}

\subsection*{You'll want to know: \kn{ವಿಮೆ}}
\textbf{\kn{ವಿಮೆ}} — \emph{vime} — \textquotedblleft{}insurance\textquotedblright{}.

\begin{grammarlens}[title={What you've built}]
One more word for this chapter.
\end{grammarlens}

\subsection*{Your turn}
\begin{itemize}
\raggedright
  \item \emph{Say it:} \emph{vime}
  \item \emph{Say it:} \emph{vime}, once more
  \item \emph{Say it:} say \emph{vime}
  \item \emph{Recall:} say the Kannada for an officer, then the Kannada for an expense, then say \emph{vime} again
\end{itemize}

\begin{tcolorbox}[breakable,colback=teal!4,colframe=teal!35!black,title={Before you move on}]
What does \kn{ವಿಮೆ} mean? (\textquotedblleft{}Insurance\textquotedblright{}.) Say it once more.
\end{tcolorbox}
\section[\texorpdfstring{\kn{ಸಾಲು} (sālu)}{ಸಾಲು (sālu)}]{\kn{ಸಾಲು} (sālu) — a queue, a line}

\label{lesson:KA-C232-salu}



\begin{quote}
Before the new one: say the Kannada for an expense, then the Kannada for insurance.
\end{quote}

\subsection*{You'll want to know: \kn{ಸಾಲು}}
\textbf{\kn{ಸಾಲು}} — \emph{sālu} — \textquotedblleft{}a queue, a line\textquotedblright{}.

\begin{grammarlens}[title={What you've built}]
One more word for this chapter.
\end{grammarlens}

\subsection*{Your turn}
\begin{itemize}
\raggedright
  \item \emph{Say it:} \emph{sālu}
  \item \emph{Say it:} \emph{sālu}, once more
  \item \emph{Say it:} say \emph{sālu}
  \item \emph{Recall:} say the Kannada for an expense, then the Kannada for insurance, then say \emph{sālu} again
\end{itemize}

\begin{tcolorbox}[breakable,colback=teal!4,colframe=teal!35!black,title={Before you move on}]
What does \kn{ಸಾಲು} mean? (\textquotedblleft{}A queue, a line\textquotedblright{}.) Say it once more.
\end{tcolorbox}
\section[\texorpdfstring{\kn{ಅನುಭವ} (anubhava)}{ಅನುಭವ (anubhava)}]{\kn{ಅನುಭವ} (anubhava) — experience}

\label{lesson:KA-C232-anubhava}



\begin{quote}
Before the new one: say the Kannada for insurance, then the Kannada for a queue.
\end{quote}

\subsection*{You'll want to know: \kn{ಅನುಭವ}}
\textbf{\kn{ಅನುಭವ}} — \emph{anubhava} — \textquotedblleft{}experience\textquotedblright{}.

\begin{grammarlens}[title={What you've built}]
One more word for this chapter.
\end{grammarlens}

\subsection*{Your turn}
\begin{itemize}
\raggedright
  \item \emph{Say it:} \emph{anubhava}
  \item \emph{Say it:} \emph{anubhava}, once more
  \item \emph{Say it:} say \emph{anubhava}
  \item \emph{Recall:} say the Kannada for insurance, then the Kannada for a queue, then say \emph{anubhava} again
\end{itemize}

\begin{tcolorbox}[breakable,colback=teal!4,colframe=teal!35!black,title={Before you move on}]
What does \kn{ಅನುಭವ} mean? (\textquotedblleft{}Experience\textquotedblright{}.) Say it once more.
\end{tcolorbox}
\section[\texorpdfstring{\kn{ರಾಜೀನಾಮೆ} (rājīnāme)}{ರಾಜೀನಾಮೆ (rājīnāme)}]{\kn{ರಾಜೀನಾಮೆ} (rājīnāme) — a resignation}

\label{lesson:KA-C232-rajiname}



\begin{quote}
Before the new one: say the Kannada for a queue, then the Kannada for experience.
\end{quote}

\subsection*{You'll want to know: \kn{ರಾಜೀನಾಮೆ}}
\textbf{\kn{ರಾಜೀನಾಮೆ}} — \emph{rājīnāme} — \textquotedblleft{}a resignation\textquotedblright{}.

\begin{grammarlens}[title={What you've built}]
One more word for this chapter.
\end{grammarlens}

\subsection*{Your turn}
\begin{itemize}
\raggedright
  \item \emph{Say it:} \emph{rājīnāme}
  \item \emph{Say it:} \emph{rājīnāme}, once more
  \item \emph{Say it:} say \emph{rājīnāme}
  \item \emph{Recall:} say the Kannada for a queue, then the Kannada for experience, then say \emph{rājīnāme} again
\end{itemize}

\begin{tcolorbox}[breakable,colback=teal!4,colframe=teal!35!black,title={Before you move on}]
What does \kn{ರಾಜೀನಾಮೆ} mean? (\textquotedblleft{}A resignation\textquotedblright{}.) Say it once more.
\end{tcolorbox}
